\documentclass[10pt]{article}
\usepackage[margin=20mm]{geometry}
\usepackage{amsmath,amssymb,mathtools,fontspec,unicode-math,xurl,hyperref}
\IfFontExistsTF{DejaVu Serif}{\setmainfont{DejaVu Serif}}{\setmainfont{Latin Modern Roman}}
\setmathfont{latinmodern-math.otf}
\setlength{\parindent}{0pt}\setlength{\parskip}{6pt}\sloppy
\title{Initial scientific review: M25-0045
Asymptotics for vectorial Allen–Cahn systems}
\author{}
\date{25 October 2026}
\begin{document}
\maketitle
Abstract: Theorem 1.16, read with the domain definition (1.67), admits a smooth boundary-touching subdomain for which its clearing-out conclusion is false. An explicit two-component solution satisfies every printed potential and energy assumption. The defect also identifies the unavailable compactly supported cutoff in Proposition 3.11. It is classified Major because this clearing-out guarantee is a stated tool in the core rectifiability argument. Its impact is restricted: the inspected later application is interior and survives an explicit interior-neighborhood restriction. No counterexample to the main interior compactness or stationarity theorems is established. Several tensor, scaling and measure-theoretic steps are checked separately, with unverified inputs identified.\par
\section*{Identity, version and review scope}
Fabrice Bethuel, Asymptotics for 2-dimensional vectorial Allen–Cahn systems, Acta Mathematica 234 (2025), 189–314, DOI `10.4310/ACTA.2025.v234.n2.a1`. The exact supplied publisher PDF `papers/0045/0045.pdf` is reviewed. Printed page equals one-based PDF page plus 188. Its SHA-256 is\par
> 412f3a54c30811f26b02c1dfb509155a
> a6211d90624806a5c67e32dbc3b2b2ab.\par
The line break does not belong to the hash. All manifest identities and file bindings were checked before scientific review. Index parsing is not counted as proof inspection.\par
Outcome: one Major finding, M01, concerning a false domain-general clearing-out statement; no separate supported Minor finding. This is a systematic claim-focused initial review, not certification of the entire 126-page paper. Text was actually read on PDF pages 1–30, 48–54, 64–66, 69–89, 109–117 and 121–124 (74 distinct pages). Twenty-five source pages were rendered and visually checked, including every decisive page for M01. The full extraction is merely a navigation artifact.\par
Equation reference: eq:pde. The standing assumptions are a smooth bounded \(\Omega\subset\mathbb R^2\), smooth nonnegative \(V:\mathbb R^k\to\mathbb R\), finitely many zeroes with positive definite Hessians, coercivity and the radial growth condition (H3), together with solutions \[
 \Delta u_\varepsilon=\varepsilon^{-2}\nabla V(u_\varepsilon),\qquad
 E_\varepsilon(u_\varepsilon)=\int_\Omega
 \left(\frac{\varepsilon}{2}|\nabla u_\varepsilon|^2+
 \frac{V(u_\varepsilon)}{\varepsilon}\right)\le M_0 .\] The limits \(\nu_\star\) and \(\zeta_\star\) refer respectively to full energy and potential energy. The main results concern interior concentration, rectifiability, local straight segments away from an \(\mathcal H^1\)-null exceptional set, and stationarity of the potential-energy varifold. No boundary condition is imposed in these hypotheses. Remark 1.4 explicitly treats Neumann boundary orthogonality as an additional question.\par
\section*{Problem 1: the clearing-out statement needs an interior restriction}
\subsection*{Source claim and failing step}
Equation reference: eq:claimed. Theorem 1.16, printed p. 211/PDF 23, asserts, for an open \(U\subset\Omega\) and \(\delta>0\), \[
 \nu_\star(V_\delta)=0\quad\Longrightarrow\quad
 \nu_\star(\overline U)=0 ,
 \qquad
 U_\delta=\{x\in\Omega:\operatorname{dist}(x,U)\le\delta\},\
 V_\delta=U_\delta\setminus U .\] The overline on \(U\) in the conclusion was confirmed in the page image; text extraction suppresses it. The preceding prose allows general regular subdomains and does not require \(\overline U\subset\Omega\) or a positive distance to the exterior boundary.\par
Section 7.3, printed p. 273/PDF 85, invokes Proposition 4.15. The latter, printed p. 252/PDF 64, writes \(U_\delta\subset\Omega\), but this is already automatic for the clipped definition (1.67). Its proof uses Proposition 3.11, printed p. 242/PDF 54, where a smooth cutoff equals one on \(U_{\delta/2}\) and vanishes outside \(U_\delta\), and the test field \(X(z)=z\chi_\delta(z)\) is inserted in the compact-support identity (3.19). Such a cutoff is unavailable when \(U_{\delta/2}\) reaches \(\partial\Omega\). This is a boundary-support obstruction, not a numerical approximation error.\par
\subsection*{Complete counterexample}
Take \[\Omega=\{x_1^2+x_2^2<1\},\quad k=2,\quad
 V(a,b)=\frac{(1-a^2)^2}{4}+\frac{b^2}{2},\qquad
 u_\varepsilon(x)=\left(\tanh\frac{x_1}{\sqrt2\,\varepsilon},0\right).\] This is genuinely a map to the allowed two-component target, even though one component is constant.\par
Potential assumptions. The zero set is \(\{(-1,0),(1,0)\}\), so (H1) holds with \(q=2\). The Hessian at either zero is \(\operatorname{diag}(2,1)\), establishing (H2). The polynomial is smooth, nonnegative and tends to infinity as \(|(a,b)|\to\infty\). Moreover, \[\begin{aligned}
 (a,b)\cdot\nabla V&=a^4-a^2+b^2\\
 &=(a^2-\tfrac34)^2+\tfrac12(a^2+b^2)+\tfrac12b^2-\tfrac9{16}\\
 &\ge \tfrac14(a^2+b^2)\quad\text{if }|(a,b)|>2.
\end{aligned}\] Thus (H3) holds with \(R_\infty=2\) and \(\alpha_\infty=1/4\).\par
Equation and uniform energy. Writing \(w=\tanh(x_1/(\sqrt2\varepsilon))\) gives \[w_{x_1x_1}=\varepsilon^{-2}(w^3-w),\qquad
 \nabla V(w,0)=(w^3-w,0).\] Consequently [eq:pde] (\#eq:pde) holds for every \(0<\varepsilon\le1\). Its energy density is \[e_\varepsilon(u_\varepsilon)
 =\frac1{2\varepsilon}\operatorname{sech}^4
       \left(\frac{x_1}{\sqrt2\varepsilon}\right).\] The integral in the \(x_1\) direction is \[c=\frac{\sqrt2}{2}\int_{\mathbb R}\operatorname{sech}^4 s\,ds
   =\frac{2\sqrt2}{3},\] because \(t=\tanh s\) changes the integral to \(\int_{-1}^1(1-t^2)\,dt=4/3\). Every vertical slice of the unit disk has length at most two. Therefore \(E_\varepsilon\le2c=4\sqrt2/3<2\) for every allowed \(\varepsilon\); choose \(M_0=2\).\par
Actual limiting measure. Changing variables \(x_1=\sqrt2\varepsilon s\) in the integral against any compactly supported continuous test function, then using the integrable dominating function \(\operatorname{sech}^4s\), proves \[\nu_\star=c\,\mathcal H^1\mathbin{\rule[0.3ex]{0.07ex}{1.1ex}\rule{0.7ex}{0.07ex}}L,\qquad
 L=\{(0,x_2):-1<x_2<1\}.\] The same calculation gives \(\zeta_\star=(c/2)\mathcal H^1\mathbin{\rule[0.3ex]{0.07ex}{1.1ex}\rule{0.7ex}{0.07ex}}L\). These limits hold along every sequence \(\varepsilon\to0\); no subsequence choice can remove the counterexample.\par
Subdomain and contradictory conclusion. Set \[U=\{4x_1^2+x_2^2<1\},\qquad \delta=\tfrac14.\] This is a bounded, connected, smooth, proper open subset of \(\Omega\). Its closure touches \(\partial\Omega\) only at \((0,\pm1)\), and \(L\subset U\). By the printed definition, \(V_\delta\subset\Omega\setminus U\), hence \[\boxed{\nu_\star(V_\delta)=0,\qquad
        \nu_\star(\overline U\cap\Omega)=2c=\frac{4\sqrt2}{3}>0.}\] If \(\nu_\star\) is extended by zero outside \(\Omega\), the second value is exactly \(\nu_\star(\overline U)\). If all sets are interpreted relatively in \(\Omega\), the value is unchanged. Thus both conventional interpretations contradict [eq:claimed] (\#eq:claimed). This example does not use equality \(U=\Omega\), an irregular subdomain, a degenerate well, unbounded energy, or a boundary condition absent from the theorem.\par
For completeness, the corresponding finite-\(\varepsilon\) exterior energy tends to zero as well. For almost every \(x_2\in(-1,1)\), the ellipse contains a positive-width interval around \(x_1=0\). The one-dimensional energy mass outside that interval tends to zero and is bounded by \(c\). Dominated convergence in \(x_2\) proves \(\int_{\Omega\setminus U}e_\varepsilon\to0\). Meanwhile the potential energy on \(U\) tends to \(c>0\). Thus Proposition 3.11’s estimate also fails in this boundary-touching range for any fixed finite constant independent of \(\varepsilon\).\par
\subsection*{Alternative readings, impact and severity}
The most favorable reading is that every subdomain in this argument was intended to be relatively compact, or that \(U_\delta\) was intended to be an ambient neighborhood contained in \(\Omega\). Either reading excludes the example. Neither is stated in Theorem 1.16, and the explicit \(x\in\Omega\) in (1.67) makes the displayed inclusion in Proposition 4.15 insufficient by itself. The correction must state an interior separation, not just repeat that inclusion.\par
The dependency is explicit: Theorem 1.16 is used in Proposition 7.4 to eliminate separated interior components; Proposition 1.17 and the rectifiability part of Theorem 1.2 then use this connectedness. A false theorem and unavailable test field on this core proof route warrant Major under the assigned calibration. Repair difficulty is not the severity criterion.\par
There is decisive counterevidence against a broader adverse conclusion. In Section 7.4, printed pp. 274–276/PDF 86–88, the application is constructed inside \(D^2(x_0,r)\) with \(D^2(x_0,2r)\subset\Omega\). The separated component and its small thickening are compactly contained in that interior disk. They satisfy the additional restriction below. The present finding therefore does not refute Theorems 1.2, 1.3, 1.5, 1.8 or 1.9; the example itself has a straight stationary interior interface.\par
\section*{A precise restricted repair}
A defensible corrected statement is: assume \(U\) is bounded and open and \[K_\delta=\{x\in\mathbb R^2:\operatorname{dist}(x,\overline U)\le\delta\}
 \Subset\Omega .\] Define \(V_\delta=K_\delta\setminus U\). Under the paper’s PDE and energy hypotheses, \(\nu_\star(V_\delta)=0\) implies \(\nu_\star(\overline U)=0\).\par
Choose \(0<\delta_0<\min\{\delta,1\}\). Then \(K_{\delta_0}\Subset\Omega\) and \(V_{\delta_0}=K_{\delta_0}\setminus U\subset V_\delta\), so \(\nu_\star(V_{\delta_0})=0\). The cutoff in Proposition 3.11 can now be chosen with compact support in \(\Omega\): take it equal to one on \(K_{\delta_0/2}\) and supported in the interior of \(K_{\delta_0}\). The stress identity gives the potential bound because derivatives of the cutoff occur only in the exterior band, where the energy controls the tensor. This permits the use of Proposition 4.15 in its interior setting and with its explicit width condition \(\delta_0<1\).\par
The passage to the measure limit can also be stated rigorously. \(V_{\delta_0}\) is compact and closed in \(\Omega\), so weak convergence and its zero limiting mass give \(\limsup\nu_{\varepsilon_n}(V_{\delta_0})\le0\). For sufficiently large \(n\) the smallness threshold of Proposition 4.15 holds, and \[\nu_{\varepsilon_n}(U_{\delta_0/4})
 \le C\{\nu_{\varepsilon_n}(V_{\delta_0})+\varepsilon_nM_0\}\longrightarrow0.\] An open neighborhood of \(\overline U\) contained in \(U_{\delta_0/4}\) then has zero limiting mass, proving the corrected conclusion. If only \(\overline U\Subset\Omega\) is assumed, choose the smaller width below \(\delta\), one, and the distance to \(\partial\Omega\); its exterior band is still a subset of the original band.\par
This verifies the support and limiting steps of the restricted repair, conditional on the earlier interior energy estimates used in Proposition 4.15. It is not an independent reconstruction of all level-set estimates in Section 4. The original boundary-general claim is not restored, and no replacement boundary theorem is proposed.\par
\section*{Other mathematical checks and their limits}
\subsection*{Scaling and the fast-decay recurrence}
Equations (1.52)–(1.61), printed pp. 208–209, have the correct two-dimensional scaling: \(\widetilde\varepsilon=\varepsilon/r\), \(\widetilde u(x)=u(rx+x_0)\) and \(E_{\widetilde\varepsilon}(\widetilde u,D)=E_\varepsilon(u,D_r)/r\). The scaled decay has terms \(r^{-1/2}E^{3/2}\) and \((\varepsilon/r)E\).\par
For the recurrence in Section 6.2, printed pp. 261–265, putting \(A_n=-\log E_n\) gives \(A_{n+1}\ge(3/2)A_n-f_n\), with \(f_n=(n/2)\log2+\log K_{\rm dec}\ge0\). Direct induction gives the exact lower bound \[A_n\ge(3/2)^n\left(A_0-\sum_{j=0}^{n-1}(2/3)^{j+1}f_j\right).\] The infinite sum is \(2\log2+2\log K_{\rm dec}\). Choosing \(A_0\) larger by one yields the double-exponential energy estimate and the logarithmic stopping-radius bound. The printed summation to \(n\) is weaker in this application precisely because \(f_n\ge0\); that endpoint cannot be used without qualification for arbitrary signed forcing. This is an index-level correction, not an additional substantive finding. The \(r\ge\varepsilon\) stopping condition is essential. The complete level-set derivation of Proposition 1.12 and every explicit constant/radius in Sections 5–6 were not certified.\par
\subsection*{Stress tensor, potential monotonicity and singular mass}
The compact-support identity in Lemma 3.5 follows by multiplying the PDE by \(X\cdot\nabla u\) and integrating by parts. The Hopf form on printed p. 237 and the computation on p. 241 agree with \[\frac{d}{dr}\frac{\zeta_\varepsilon(D_r)}r
 =\frac1{4r}\int_{\partial D_r}
 \left(\frac{2V(u_\varepsilon)}{\varepsilon}
       +\varepsilon|u_r|^2-\varepsilon|u_\tau|^2\right).\] No positivity at fixed \(\varepsilon\) is assumed. Once Proposition 1.20 supplies the absolutely continuous discrepancy relation, rotation of the quadratic tensor gives \(N^{\rm ac}_{x_0,\star}=4\sin^2(\gamma_\star-\theta)\zeta_\star^{\rm ac}\ge0\), as in Lemma 9.1.\par
The singular-part argument in Lemmas 9.4–9.6 was checked as an argument about measures, not pointwise derivatives. With radial projection \(\Pi(x)=|x-x_0|\), the relevant exceptional set of radii is \(\Pi(E_\star)\), which is Lebesgue-null since \(\Pi\) is Lipschitz. On its complement the pushforward of \(N\) is positive. On that set, the identity \[\Pi_\#N=4\left(\Pi_\#\zeta_\star-\frac{Z(r)}r\,dr\right)\] is positive because the absolutely continuous \(Z(r)dr/r\) term vanishes there. This establishes positivity of the distributional derivative of \(Z(r)/r\). It does not assume away the possible singular mass at the start. The displayed inverse-image symbol and dropped positive factors in the surrounding text are not needed for this verified argument and are not counted as separate scientific issues.\par
Corollary 9.8’s abbreviated appeal to a finite density at \(\lambda\)-almost every point would not, in isolation, exclude a singular measure. The available stronger input is Lemma 9.7 for every interior center together with the locally uniform bound \(\zeta_\star(D_r(x))\le C r\) from monotonicity. It gives \(\nu_\star(D_s(x))\le C' s\) uniformly on compact subsets. Covering an \(\mathcal H^1\)-null set by small balls proves zero \(\nu_\star\) mass. Thus this step is supportable without the invalid general inference from almost-everywhere finiteness alone. No optimal numerical constant in (9.21) is asserted by this review.\par
The derivation of Proposition 1.20 from the projected-moment analysis in Section 8 is a substantive unverified input here. The conditional checks above do not certify it.\par
\subsection*{Stationarity and the distinction between the two energies}
In a tangent-normal frame, write the gradient-density matrix as \(\operatorname{diag}(a,a+2\Theta)\) using Theorem 1.9. Then \(\nu=a+2\Theta\), and the limiting stress is \[\nu I-\mu=2\Theta\,t\otimes t .\] Consequently its divergence-free identity is exactly \(\int \operatorname{div}_{T_xS_\star}X\,d\zeta_\star=0\). This independently checks the sign and factor in Section 11, printed pp. 311–312. Tangential energy \(a\) need not vanish. Therefore stationarity of \(\zeta_\star\) does not, by itself, assert \(\nu_\star=2\zeta_\star\). No unsupported energy-equipartition substitution is needed.\par
\section*{Overall disposition and remaining work}
The proof route supports an interior interpretation of the clearing-out tool, and the inspected connectedness use satisfies that interpretation. The unrestricted printed formulation needs correction. Neither arbitrary boundary behavior nor higher-dimensional or parabolic extensions follow from this review.\par
Item | Disposition\par
M01 | Major: boundary-touching case falsifies Theorem 1.16 as stated and the unrestricted supporting estimate.\par
Repair | Interior restriction and limiting step verified, conditional on the earlier interior energy estimates.\par
Main results | No counterexample established; inspected local use survives the restriction.\par
Other checks | Scaling, recurrence structure, tensor algebra and final measure arguments checked within stated dependencies.\par
Adoption/acceptance | Author adoption not requested; independent human scientific acceptance pending.\par
Remaining work | Full level-set estimates, Section 8 projected moments, all Section 10 steps, quantitative constants, and cited geometric-measure-theory foundations.\par
No separate supplement or computational repository is included in the acquired manifest. No author code was run; its absence is not evidence of correctness or a defect. The pseudo-profile construction cited as work in preparation was not reconstructed. The complete paper and all external dependencies have not been verified.\par
\section*{Supporting checks}
The source-page images actually inspected are PDF pages 1, 4, 6, 8, 9, 18, 21, 23, 49, 53, 54, 64, 73, 74, 75, 85, 86, 87, 88, 111, 114, 115, 116, 123 and 124. The full extracted text, source PDF information and actual I/O events are retained alongside them. An output-size truncation during a combined read was resolved by rereading PDF pages 48–50; those pages were not credited from truncated output.\par
The reviewer-authored `independent\_checks.py`, in the evidence directory, was actually executed with existing Windows Python. It checks the rational mass integral, all four polynomial coefficients in the corrected (H3) identity, the positive radial-growth margin at radius two, and twenty exact recurrence steps using only the standard library. The initial report omitted \(b^2/2\) from the identity and its original script checked only two coefficients; focused acceptance corrected both and executed the updated check. The JSON result records the analytic counterexample parameters and exact values. Finite recurrence checks corroborate arithmetic; induction, PDE satisfaction and weak convergence are proved above, not inferred from those checks. This is not empirical reproduction.\par
Version checks consulted arXiv `2003.10189`, whose inspected submission history lists v1, 23 March 2020. It was not substituted for the supplied 2025 publisher PDF. Targeted correction searches did not identify a correction for this theorem; this limited negative result does not prove no correction exists or establish novelty. The publisher landing-page open returned HTTP 403; the supplied PDF remained available, and its identity was visually verified. No claim of reading the blocked landing page is made.\par
The initial native build is preserved. The current exact TeX is compiled with the existing authorized Tectonic runtime; its command, dependency hashes, actual result and rendered-page checks are recorded in the current build record. No author software was executed.\par
Human acceptance remains pending.\par
9 F. Bethuel, Asymptotics for 2-dimensional vectorial Allen–Cahn systems, Acta Math. 234 (2025), 189–314. DOI: `10.4310/ACTA.2025.v234.n2.a1`. F. Bethuel, Asymptotics for two-dimensional vectorial Allen–Cahn systems, arXiv:2003.10189v1, 23 March 2020. https://arxiv.org/abs/2003.10189. Submission metadata inspected; no claim of full-version equivalence.\par
\end{document}
